Spectral analysis begins after the graph has been defined, not before.  This ordering sounds
obvious, yet protein sequence analysis offers many matrices with plausible eigenvectors: an MSA
covariance, a Potts norm matrix, an attention route, a dynamic response shadow, a signed coupling
graph, or a connection operator acting on categorical fibers.  Their spectra live in different
spaces and answer different questions.

The first decision is what has been compressed.  A scalar Laplacian describes smoothness of
functions over residue positions relative to nonnegative edge weights.  A signed Laplacian retains
one oriented scalar polarity.  A connection Laplacian transports vectors between local
categorical spaces and can test whether substitution directions remain coherent around the graph.
No one spectrum contains the others.

Directionality creates an earlier branch.  A row-softmax route is naturally a Markov operator.
Its Euclidean asymmetry need not imply irreversible circulation: a symmetric score kernel can
produce an entrywise asymmetric but reversible route under its stationary measure.  Conversely, a
doubly stochastic route can remain nonreciprocal.  Reversibility, detailed-balance defect, and
directed singular modes should be measured before projecting the route to an undirected shadow.

Context dependence adds another noncommutativity.  The spectrum of a mean graph, the mean spectrum
of context-specific graphs, and the distribution of spectral projectors are different objects.
Near degeneracy makes individual eigenvectors unstable even when their invariant subspace is
stable.  Spectral claims should therefore be phrased at the level of the appropriate projector and
calibrated against replicate null graphs rather than inferred from rank alone.

The guiding question of the chapter is not ``what does the leading eigenvector mean?'' but ``which
operator, on which space, relative to which null, produced this eigenvector?''  Only then can a
spectral mode be connected responsibly to conservation, sectors, communication, allostery, or
another biological interpretation.
